\section{Background: why MPI worked}
\label{sec:background}

MPI is often described as a library of communication functions. It is more
useful, for our purposes, to read it as a set of decisions about
\emph{what to standardise and what to leave alone}. This section extracts the
decisions we inherit. Readers who know MPI can skip to
\Cref{sec:mismatch}; readers who know agent systems and not MPI will find
that most of what follows describes problems they have already met.

\subsection{The standard is not an implementation}

The Forum standardised an interface and refused to standardise
behaviour below it~\cite{mpi-forum-1994}. \mpiop{Bcast} says what every
rank ends up holding; it says nothing about the schedule. MPICH and Open MPI
between them implement at least five broadcast algorithms and choose among
them at runtime by message size, communicator size and
topology~\cite{Thakur2005,Chan2007}. Applications
written in 1996 got faster in 2016 without being touched.

This is the decision that separates a protocol from a framework, and it is
the one the agent ecosystem has not yet made. A harness written against
AutoGen's group chat encodes a scheduling policy in the application; a
harness written against \mpiop{Bcast} encodes an intent.

\subsection{Communicators, contexts, and safe composition}
\label{sec:bg-contexts}

A communicator pairs an ordered group of ranks with an opaque
\emph{communication context}. Two properties follow. First, a rank is
addressed by its index \emph{within the communicator}, so a subgroup can be
written as if it were the whole world. Second, and much more importantly, a
message sent on one communicator can never be matched by a receive on
another.

The second property is what made MPI libraries possible. Zipcode's designers
had identified the problem directly: a library that posts a wildcard receive
can consume a message the application sent to itself, and no discipline in
tag allocation fixes it once two independently written libraries are
involved~\cite{skjellum1994zipcode}. \mpiop{Comm\_dup} makes the fix a
one-liner: a library duplicates the communicator it is handed and is
thereafter immune.

Multi-agent frameworks are, almost uniformly, back in the pre-context world.
The shared conversation buffer is a flat tag space, and a sub-agent that
reads ``the last message'' has no way to know whether it was addressed to
it. This is not a hypothetical: it is one of the failure classes the MAST
taxonomy names~\cite{cemri2025mast}.

\subsection{A small orthogonal core, plus a large convenience layer}

MPI-1's point-to-point layer is, in principle, four calls. Everything
else---the other three send modes, the nonblocking forms, the persistent
requests, and every collective---could be written on top. The Forum included
them anyway, for a reason that generalises: a collective operation
\emph{states an intent} that the implementation can exploit, while the
hand-rolled equivalent states a schedule that it cannot. An \mpiop{Allreduce}
can be a recursive-doubling butterfly; a loop of sends and receives that
computes the same thing cannot be.

The Forum's discipline about what to include was real and enforced---a
feature needed a working prototype before it could be voted---and it has an
instructive cost. TCGMSG offered \code{nxtval}, an atomic shared counter for
dynamic load balancing. It had no MPI counterpart, was excluded as not being
message passing, and was subsequently re-implemented over MPI in something
like ten platform-specific ways, one of which burns an entire process as a
counter server. We note it because that primitive is
\op{Fetch\_and\_op} on a window (\Cref{sec:design-rma}), it is how an
AgentMPI harness implements a shared work queue with no orchestrator, and a
protocol for executors with order-of-magnitude variation in per-task cost
cannot afford to leave it out.

\subsection{Collective algorithms and their cost model}

Collective algorithm design is where MPI's engineering depth lies. The
standard cost model is Hockney's: a message of $n$ bytes between two ranks
costs $\alpha + n\beta$, with $\alpha$ the latency and $\beta$ the
reciprocal bandwidth~\cite{Hockney1994}; LogP and its
descendants refine it with overhead and gap
terms~\cite{Culler1993,Alexandrov1995}. In this model,
broadcast by binomial tree costs $\lceil\log_2 p\rceil(\alpha + n\beta)$
while van de Geijn's scatter-then-allgather costs roughly
$(\log_2 p + p - 1)\alpha + 2\frac{p-1}{p}n\beta$, so the tree wins for
small messages and the scatter/allgather form wins for large
ones~\cite{Thakur2005,Chan2007}. Allreduce shows the
same structure: recursive doubling costs $\log_2 p$ rounds but moves $n$
bytes in each, while Rabenseifner's reduce-scatter-plus-allgather moves
$2\frac{p-1}{p}n$ bytes in $2\log_2 p$
rounds~\cite{Rabenseifner2004}. The ring allreduce, later
rediscovered by the deep-learning community, is bandwidth-optimal and
latency-terrible: $2(p-1)$ rounds~\cite{Patarasuk2009}.

Every one of these trade-offs is a statement about the relative size of
$\alpha$ and $\beta$. \Cref{sec:cost} shows what happens when the constants
change by six orders of magnitude.

\subsection{One-sided operations}

MPI-2 added remote memory access because two-sided messaging forces the
receiver to participate in every transfer, and there are algorithms where
only the initiator knows what it needs~\cite{mpi-forum-1997-mpi20}. A rank
exposes a \emph{window}; peers put, get and accumulate into it within
synchronisation epochs---collective \mpiop{Win\_fence}, scoped
post/start/complete/wait, or passive-target lock/unlock. MPI-3 tightened the
memory model and added atomics and shared-memory
windows~\cite{dinan2016rma}.

The agent analogue is the shared scratchpad that every framework has and
none of them specifies. What MPI-2 supplies and the scratchpads do not is
epochs: a defined moment before which writes are not visible and after which
they are. Without one, a reader sees a torn state and reasons confidently
about half an answer.

\subsection{Collective I/O}

Naive parallel I/O collapses: a thousand ranks each writing a small
unaligned piece of a shared file produce a thousand conflicting requests.
ROMIO's two-phase I/O redistributes data to a small set of aggregators so
that each issues one large contiguous write~\cite{Thakur1999}. An
agent harness writing to a shared repository has the same problem with a
worse failure mode: twenty agents each committing a two-line edit produce
merge conflicts, and a merge conflict costs an agent turn.

\subsection{Fault tolerance: ULFM}

MPI's position from 1994 until today is that after a process failure the
state of MPI is undefined, with \mpiop{ERRORS\_ARE\_FATAL} the default. It
was a deliberate choice, and the standard says so in a goal statement worth
quoting because it is exactly the assumption an agent runtime cannot
make: MPI is to \emph{``assume a reliable communication interface: the user
need not cope with communication failures''}~\cite{mpi-forum-1994}. On a
machine whose mean time between failures is days and whose jobs checkpoint
anyway, that was defensible.

The consequences have proved remarkably durable. The User Level Failure
Mitigation proposal adds the minimum viable set of primitives on
top~\cite{bland2012ulfm,Bland2013a}, and after fifteen years it is still a
proposal: its first slice passed a final vote in February 2023 but sits on
\code{mpi-next}, the \mpiop{Comm\_agree} slice failed quorum in December
2023, and the \mpiop{Comm\_shrink} slice failed quorum in June
2026.\footnote{Contrary to several secondary sources, MPI-5.0 (June 2025)
contains no fault-tolerance chapter and no occurrence of \code{revoke},
\code{shrink}, \code{agree} or \code{MPI\_ERR\_PROC\_FAILED}; it is
substantially a single-feature release standardising an
ABI~\cite{mpi-forum-1994}.} That a well-specified, prototyped, actively
championed recovery model has resisted retrofitting for fifteen years is the
strongest available argument for putting one in a new protocol at the start,
which is what we do. The primitives themselves: \mpiop{Comm\_revoke} invalidates a
communicator everywhere, so that a rank blocked on a peer that will never
send gets an error instead of waiting forever; \mpiop{Comm\_shrink} builds a
new communicator over the survivors and must be an agreement, since two
ranks that disagreed about who died would build different worlds; and
\mpiop{Comm\_agree} is a fault-tolerant agreement, which the FLP result tells
us requires a failure detector~\cite{Fischer1985}. ULFM's
design principle---\emph{the library should not recover for you, it should
make recovery writable}---is one we adopt verbatim.

The earlier FT-MPI took a different line, offering a \code{REBUILD} mode
that respawned failed ranks and preserved their
indices~\cite{Fagg2000}. For MPI that was expensive, because a rank's state
is gigabytes of application memory. \Cref{sec:ft} argues that for agents it
is the right default, because a rank's state is a prompt and a checkpoint.

\subsection{Tooling: PMPI and MPI\_T}

MPI standardised a name-shifted profiling interface in 1994: every
\mpiop{X} has a weak symbol over \code{PMPI\_X}, so a tool can interpose
without touching any implementation. An entire ecosystem---mpiP, TAU,
Score-P, Vampir, Jumpshot---exists because of that one
decision~\cite{shende2006tau,knupfer2012scorep}. MPI-3 added \mpiop{T} for
reading and writing an implementation's internal control and performance
variables~\cite{islam2016mpit}. Agent frameworks have nothing comparable;
observability is per-framework and incompatible. We treat both as part of
the protocol rather than as implementation features, for the same reason the
Forum did.

\subsection{The process manager is out of scope}

Finally, MPI says almost nothing about how processes start. Launch and
wire-up belong to a process manager speaking a separate
protocol---PMI, and later PMIx, which exists precisely because exchanging
endpoint information among $10^5$ ranks through an argument list does not
work~\cite{castain2018pmix}. The separation is why the same MPI library runs
under Slurm, under a laptop, and under schedulers that did not exist when it
was written. \ampi{} draws the same line, and \Cref{sec:impl} shows what
lies on each side of it.
